\chapter*{ระเบียบความปลอดภัย จรรยาบรรณห้องปฏิบัติการ และการเขียนรายงาน}
\addcontentsline{toc}{chapter}{ระเบียบความปลอดภัย จรรยาบรรณห้องปฏิบัติการ และการเขียนรายงาน}

การปฏิบัติการทดลองในรายวิชา 4011111 ฟิสิกส์สำหรับจุลชีววิทยา มีเป้าหมายสำคัญในการปลูกฝัง \textbf{ผลลัพธ์การเรียนรู้รายวิชาด้านจรรยาบรรณและความปลอดภัย (CLO3 เชื่อมโยงกับ PLO1)} และ \textbf{ทักษะการทำงานร่วมกับผู้อื่นและการสื่อสารทางวิทยาศาสตร์ (CLO4 เชื่อมโยงกับ PLO5 และ PLO6)} นักศึกษาทุกคนต้องปฏิบัติตามระเบียบต่อไปนี้อย่างเคร่งครัด

\section*{1. กฎความปลอดภัยในห้องปฏิบัติการฟิสิกส์และจุลชีววิทยา (Laboratory Safety)}
\begin{enumerate}
    \item \textbf{การแต่งกายและอุปกรณ์ป้องกันส่วนบุคคล}\quad นักศึกษาต้องสวมเสื้อกาวน์ปฏิบัติการ (Lab Coat) ติดกระดุมเรียบร้อย สวมรองเท้าหุ้มส้นมิดชิด รวบผมยาวให้เรียบร้อย และห้ามนำอาหารหรือเครื่องดื่มเข้ามารับประทานในห้องปฏิบัติการโดยเด็ดขาด
    \item \textbf{ความปลอดภัยด้านเครื่องปั่นเหวี่ยง (Centrifuge)}\quad ต้องตรวจสอบว่าหลอดตัวอย่างที่วางตรงข้ามกันมีปริมาตรและมวลเท่ากันเสมอ ปิดฝาครอบโรเตอร์ให้สนิท และห้ามเปิดฝาเครื่องจนกว่าโรเตอร์จะหยุดหมุนสนิท ($0\text{ rpm}$)
    \item \textbf{ความปลอดภัยด้านความร้อนและหม้อนึ่งความดันไอ (Autoclave)}\quad ห้ามบิดฝาขวดแก้วจนแน่นสนิทก่อนนำเข้าหม้อนึ่ง (ต้องคลายเกลียวฝาออกครึ่งรอบเสมอ) และห้ามเปิดฝาหม้อนึ่งขณะที่ความดันเกจยังมากกว่า $0\text{ psi}$ หรืออุณหภูมิสูงกว่า $80\,^\circ\text{C}$
    \item \textbf{ความปลอดภัยด้านไฟฟ้าและรังสี UV-C}\quad เช็ดมือให้แห้งสนิทก่อนสัมผัสแหล่งจ่ายไฟของชุด Gel Electrophoresis ปิดฝาครอบนิรภัยก่อนเปิดสวิตช์ไฟทุกครั้ง และห้ามจ้องมองหลอดรังสี UV-C ด้วยตาเปล่าโดยไม่มีแผ่นกรองป้องกัน
\end{enumerate}

\section*{2. จรรยาบรรณแห่งวิชาชีพวิทยาศาสตร์ในการบันทึกข้อมูล (Data Integrity — CLO3)}
\begin{itemize}
    \item \textbf{ห้ามดัดแปลงหรือเมคตัวเลขผลการทดลอง (Zero Tolerance for Data Fabrication)}\quad นักศึกษาต้องบันทึกค่าที่อ่านได้จริงจากเครื่องมือวัดลงในตารางบันทึกผลด้วยปากกา แม้ว่าค่าที่วัดได้จะมีความคลาดเคลื่อนไปจากค่าทางทฤษฎีก็ตาม
    \item \textbf{การวิเคราะห์สาเหตุความคลาดเคลื่อนมีคุณค่ามากกว่าตัวเลขที่ตรงทฤษฎีแบบปลอมๆ}\quad ผู้สอนให้คะแนนเต็มสำหรับกลุ่มที่รายงานค่าจริงและสามารถอธิบายสาเหตุทางฟิสิกส์ของความคลาดเคลื่อนได้ เช่น การสูญเสียความร้อนสู่สิ่งแวดล้อม หรือฟองอากาศติดที่ปลายทิปไมโครปิเปต
\end{itemize}

\section*{3. โครงสร้างมาตรฐานของการเขียนรายงานปฏิบัติการรายกลุ่ม (Lab Report — CLO4)}
รายงานปฏิบัติการแต่ละสัปดาห์มีคะแนนเต็ม 10 คะแนน (รวมตลอดภาคการศึกษาคิดเป็น 20\% ของคะแนนรวมรายวิชา) โดยประกอบด้วย 5 ส่วนสั้นกระชับ ได้แก่
\begin{enumerate}
    \item \textbf{ชื่อปฏิบัติการ วัตถุประสงค์ และรายชื่อสมาชิกกลุ่มพร้อมระบุหน้าที่รับผิดชอบ} (1 คะแนน)
    \item \textbf{ตารางบันทึกผลการทดลองจริงที่ระบุหน่วยวัด SI และเลขนัยสำคัญถูกต้อง} (3 คะแนน)
    \item \textbf{การแสดงวิธีคำนวณตัวอย่างและการเขียนกราฟที่มีชื่อแกนและหน่วยครบถ้วน} (3 คะแนน)
    \item \textbf{การสรุปและอภิปรายผลการทดลองเชื่อมโยงสู่งานจุลชีววิทยาและวิเคราะห์ความคลาดเคลื่อน} (2 คะแนน)
    \item \textbf{การตอบคำถามท้ายปฏิบัติการครบทุกข้อ} (1 คะแนน)
\end{enumerate}
